\begin{mvtable}
\phantomsection\label{tab:5.8}\mvtabcap{表 5.8\quad 威斯康星州麦迪逊警署的 5 类加班小时数}
\begin{tabular}{ccccc}
\toprule
\shortstack{$x_1$ 法定加班\\出勤时数} & \shortstack{$x_2$ 特殊事件\\处理加班时数} & \shortstack{$x_3$ 处理遗留\\人员加班时数} & \shortstack{$x_4$ 允许补偿\\的加班时数} & \shortstack{$x_5$ 会议时数}\\
\midrule
3\,387 & 2\,200 & 1\,181 & 14\,861 & 236\\
3\,109 & 875 & 3\,532 & 11\,367 & 310\\
2\,670 & 957 & 2\,502 & 13\,329 & 1\,182\\
3\,125 & 1\,758 & 4\,510 & 12\,328 & 1\,208\\
3\,469 & 868 & 3\,032 & 12\,847 & 1\,385\\
3\,120 & 398 & 2\,130 & 13\,979 & 1\,053\\
3\,671 & 1\,603 & 1\,982 & 13\,528 & 1\,046\\
4\,531 & 523 & 4\,675 & 12\,699 & 1\,100\\
3\,678 & 2\,034 & 2\,354 & 13\,534 & 1\,349\\
3\,238 & 1\,136 & 4\,606 & 11\,609 & 1\,150\\
3\,135 & 5\,326 & 3\,044 & 14\,189 & 1\,216\\
5\,217 & 1\,658 & 3\,340 & 15\,052 & 660\\
3\,728 & 1\,945 & 2\,111 & 12\,236 & 299\\
3\,506 & 344 & 1\,291 & 15\,482 & 206\\
3\,824 & 807 & 1\,365 & 14\,900 & 239\\
3\,516 & 1\,223 & 1\,175 & 15\,078 & 161\\
\bottomrule
\end{tabular}
\end{mvtable}

我们来考察法定加班出勤时数的稳定性,由计算机算得 $\bar{x}_1=3\,558$,由于只将单个数值进行标绘,故 $\bar{x}_1$ 与 $\bar{x}$ 相同. 此外,样本标准差为 $\sqrt{s_{11}}=607$. 于是控制限为
\[
\mathrm{UCL}=\bar{x}_1+3\sqrt{s_{11}}=3\,558+3\times607=5\,379
\]
\[
\mathrm{LCL}=\bar{x}_1-3\sqrt{s_{11}}=3\,558-3\times607=1\,737
\]
数据、中心线及控制限三者合在一起,构成\hyperref[fig:5.5]{图 5.5} 中的 $\bar{X}$ 控制图.

\begin{mvfigure}
\begin{tikzpicture}[x=0.52cm,y=1.15cm]
  \draw[->] (0,0) -- (17,0) node[right,font=\small]{观测值序号};
  \draw[->] (0,0) -- (0,4.6) node[above,font=\small]{$x_1$ 值};
  \foreach \y in {0,1,2,3,4} \draw (0,\y) -- (-0.18,\y);
  \node[left,font=\scriptsize] at (0,0) {$1\,500$};
  \node[left,font=\scriptsize] at (0,1) {$2\,500$};
  \node[left,font=\scriptsize] at (0,2) {$3\,500$};
  \node[left,font=\scriptsize] at (0,3) {$4\,500$};
  \node[left,font=\scriptsize] at (0,4) {$5\,500$};
  \draw[dashed] (0,3.879) -- (16.5,3.879) node[right,font=\scriptsize]{UCL=5\,379};
  \draw[dashed] (0,0.237) -- (16.5,0.237) node[right,font=\scriptsize]{LCL=1\,737};
  \draw[dashed] (0,2.058) -- (16.5,2.058) node[right,font=\scriptsize]{$\bar{x}_1$=3\,558};
  \draw (1,1.887)--(2,1.609)--(3,1.170)--(4,1.625)--(5,1.969)--(6,1.620)--(7,2.171)--(8,3.031)--(9,2.178)--(10,1.738)--(11,1.635)--(12,3.717)--(13,2.228)--(14,2.006)--(15,2.324)--(16,2.016);
  \foreach \x/\y in {1/1.887,2/1.609,3/1.170,4/1.625,5/1.969,6/1.620,7/2.171,8/3.031,9/2.178,10/1.738,11/1.635,12/3.717,13/2.228,14/2.006,15/2.324,16/2.016}
    \fill (\x,\y) circle (1.3pt);
\end{tikzpicture}
\phantomsection\label{fig:5.5}\mvcap{图 5.5\quad $x_1$=法定加班出勤时数的 $\bar{X}$ 控制图}
\end{mvfigure}

法定加班出勤时数在数据收集期内是稳定的,加班时数的波动是由一般原因引起的,并无迹象表明存在引起波动的特殊原因.\hfill$\blacksquare$

在考虑多于一个重要特征的情形下,需要用多元的方法来监控过程的稳定性. 这种方法能考虑诸特征之间的相关性,且能控制对波动的特殊原因发出错误信号的总体概率. 变量间的高度相关使得评估一个总的误差率成为不可能,而这点却可以通过大量一元控制图来
